% EE147 Exam 2 cheat sheet. Kaushik Vada. One sheet, front and back.
\documentclass[10pt]{extarticle}
\usepackage[landscape,letterpaper,margin=0.28in]{geometry}
\usepackage{amsmath,amssymb}
\usepackage{multicol}
\usepackage{listings}
\usepackage{xcolor}
\usepackage{enumitem}
\usepackage[T1]{fontenc}

\setlength{\parindent}{0pt}
\setlength{\parskip}{1pt}
\setlength{\columnsep}{8pt}
\setlength{\columnseprule}{0.3pt}
\pagestyle{empty}
\setlist{nosep,leftmargin=8pt}

\definecolor{hdr}{rgb}{0.12,0.25,0.55}
\definecolor{codebg}{rgb}{0.95,0.95,0.95}
\definecolor{kw}{rgb}{0,0,0.55}

\newcommand{\sect}[1]{\vspace{3pt}\par\colorbox{hdr}{\textcolor{white}{\textbf{\small #1}}}\par\vspace{2pt}}
\newcommand{\sub}[1]{\vspace{2pt}\par\textbf{\footnotesize\textcolor{hdr}{#1}}\par}

\lstset{
  basicstyle=\fontsize{7.5}{8.6}\selectfont\ttfamily,
  backgroundcolor=\color{codebg},
  keywordstyle=\color{kw}\bfseries,
  commentstyle=\color{gray},
  language=C++, breaklines=true,
  aboveskip=1.5pt, belowskip=1.5pt,
  frame=none, xleftmargin=2pt,
  morekeywords={__global__,__device__,__shared__,__syncthreads,cudaMallocManaged,
  cudaMemPrefetchAsync,cudaMemAdvise,atomicAdd,cudaDeviceSynchronize,
  cudaMemcpyAsync,cudaMemcpyPeerAsync,cudaSetDevice,cudaDeviceEnablePeerAccess,
  cudaDeviceCanAccessPeer,blockIdx,threadIdx,blockDim,gridDim,cudaMalloc,cudaFree,cudaMemcpy}
}

\begin{document}
\footnotesize

\begin{center}
\textbf{EE147 Exam 2 Cheat Sheet (Lectures 8 to 14)} \textbullet{} Kaushik Vada
\end{center}
\vspace{-4pt}

\begin{multicols*}{4}

\sect{L8: Unified Memory}
\sub{Core idea}
One pointer via \texttt{cudaMallocManaged(\&p,N)} usable from CPU and GPU; runtime migrates data. Replaces both \texttt{malloc} and \texttt{cudaMalloc}. UM is for \textbf{programmability}, not speed; tuned manual \texttt{cudaMemcpy} usually wins.

\sub{Two regimes (memorize)}
\textbf{Pre-Pascal (Kepler/Maxwell):} whole allocation migrates en masse at kernel launch; \texttt{cudaDeviceSynchronize()} brings it back. No oversubscription, no concurrent CPU/GPU access (CPU touch during kernel = \textbf{segfault}), no on-demand paging.\\
\textbf{Pascal+:} true demand paging. GPU touches non-resident page $\Rightarrow$ page fault $\Rightarrow$ driver migrates that page $\Rightarrow$ access replays. Allows oversubscription ($>$GPU mem), concurrent CPU/GPU access, system-wide atomics. Many faults = slow.

\sub{Pattern}
\begin{lstlisting}
char *d; cudaMallocManaged(&d,N);
fread(d,1,N,fp);
qsort<<<...>>>(d,N,1,cmp);
cudaDeviceSynchronize(); // before CPU use
use_data(d); cudaFree(d);
\end{lstlisting}

\sub{Tuning APIs (Pascal+)}
\texttt{cudaMemPrefetchAsync(p,len,dev,stream)} \textbf{moves} data ahead of time (\texttt{dev} may be \texttt{cudaCpuDeviceId}).\\
\texttt{cudaMemAdvise(p,n,hint,dev)} hints \textbf{never move data}:
\begin{itemize}
\item \texttt{SetReadMostly}: local read-only copy per processor; a write collapses all copies.
\item \texttt{SetPreferredLocation}: suggests home; P2P map if possible, else migrate on fault.
\item \texttt{SetAccessedBy}: direct P2P mapping for that processor, no faults, no movement.
\end{itemize}
\texttt{atomicAdd\_system}: atomic across CPU + all GPUs (Pascal+; direct on NVLink, SW-assisted on PCIe).

\sub{UM shines for}
Deep copies (struct w/ pointers), linked lists, C++ classes (overload new/delete), sparse/unpredictable access.

\sect{L9: MatMul \& Tiling}
\sub{Basic kernel}
\begin{lstlisting}
int Row=blockIdx.y*blockDim.y+threadIdx.y;
int Col=blockIdx.x*blockDim.x+threadIdx.x;
if(Row<W && Col<W){ float Pv=0;
  for(int k=0;k<W;++k)
    Pv+=M[Row*W+k]*N[k*W+Col];
  P[Row*W+Col]=Pv; }
\end{lstlisting}
2 global loads per 2 FLOPs: ratio 1, memory bound.

\sub{Tiling, 6 steps}
1) identify shared tile 2) cooperative load to \texttt{\_\_shared\_\_} 3) \texttt{\_\_syncthreads()} 4) consume from shared 5) \texttt{\_\_syncthreads()} 6) next tile.

\sub{Tiled SGEMM}
\begin{lstlisting}
__shared__ float Mds[T][T], Nds[T][T];
int bx=blockIdx.x, by=blockIdx.y;
int tx=threadIdx.x, ty=threadIdx.y;
int Row=by*T+ty, Col=bx*T+tx;
float Pv=0;
for(int p=0;p<W/T;++p){
  Mds[ty][tx]=M[Row*W + p*T+tx];
  Nds[ty][tx]=N[(p*T+ty)*W + Col];
  __syncthreads();
  for(int k=0;k<T;++k)
    Pv+=Mds[ty][k]*Nds[k][tx];
  __syncthreads();
}
P[Row*W+Col]=Pv;
\end{lstlisting}

\sub{Tile size math (exam favorite)}
\textbf{T=16:} 256 thr/blk; per phase 512 loads, 8192 mul/add $\Rightarrow$ \textbf{16 FLOPs/load}. Shared = $2{\cdot}16^2{\cdot}4$B = \textbf{2KB} $\Rightarrow$ 8 blocks in 16KB SM.\\
\textbf{T=32:} 1024 thr/blk; 2048 loads, 65536 ops $\Rightarrow$ \textbf{32 FLOPs/load}. Shared = \textbf{8KB} $\Rightarrow$ only 2 blocks/SM.\\
General: FLOPs per global load $= T$. Shared/blk $= 2T^2{\cdot}4$ bytes.

\sub{Boundary tests (3 distinct!)}
Load M: \texttt{Row<H \&\& p*T+tx<W} else store 0.\\
Load N: \texttt{p*T+ty<H \&\& Col<W} else store 0.\\
Write P: \texttt{Row<H \&\& Col<W}.\\
Zeros make mul-add a no-op; threads with no valid P still help load tiles.

\sect{L10: Memory Coalescing}
\sub{DRAM mechanics}
Cell = capacitor + transistor; random access slow. \textbf{Burst:} one row read fills buffer, streamed in N transfers at interface speed (DDR2/3 core = 1/4 or 1/8 interface speed; buffer 8x interface width). \textbf{Banks} overlap to hide dead time. HBM2/3 = stacked dies + TSVs (A100 1555 GB/s, H100 3 TB/s).

\sub{The rule}
Coalesced iff index $=$ \texttt{A[(expr indep of threadIdx.x) + threadIdx.x]}; consecutive threads hit consecutive addresses in one burst section.

\sub{MatMul example}
Row-major: \texttt{M[i][j]=M[i*W+j]}.\\
\texttt{N[k*W+Col]}, \texttt{Col=bx*T+tx}: consecutive in tx $\Rightarrow$ \textbf{coalesced}.\\
\texttt{M[Row*W+k]}: address varies with ty not tx; column-walk patterns like \texttt{A[tid*W]} are \textbf{not} coalesced.\\
\textbf{Tiling fixes coalescing:} load tile coalesced into shared, then read shared freely (no coalescing rules in shared).

\sect{L11: Histogram / Atomics}
\sub{Partitioning}
\textbf{Sectioned} (contiguous chunk per thread): adjacent threads far apart $\Rightarrow$ NOT coalesced. \textbf{Interleaved} (stride = total threads): adjacent threads adjacent memory $\Rightarrow$ coalesced. Use interleaved on GPU.

\sub{Data race}
\texttt{Old=Mem[x]; New=Old+1; Mem[x]=New} with 2 threads from 0 ends as \textbf{1 or 2}. Need atomic RMW: hardware serializes per address.

\sub{Atomics}
\texttt{int atomicAdd(int*,int)} returns \textbf{OLD} value. Variants: Sub, Min, Max, Inc, Dec, Exch, CAS, And, Or, Xor. Scopes: \texttt{\_system} (CPU+GPUs), default (device), \texttt{\_block}.

\sub{Interleaved kernel}
\begin{lstlisting}
int i=threadIdx.x+blockIdx.x*blockDim.x;
int stride=blockDim.x*gridDim.x;
while(i<size){
  int pos=buf[i]-'a';
  if(pos>=0 && pos<26)
    atomicAdd(&histo[pos/4],1);
  i+=stride; }
\end{lstlisting}

\sub{Atomic cost}
Each RMW = read latency + write latency, location locked meanwhile. Contended bin: throughput $\approx 1/\text{latency} \ll$ bandwidth. \textbf{DRAM atomic $\sim$1000+ cyc; L2 $\sim$1/10 of that; shared mem very short, $\sim$100x faster than DRAM.}

\sub{Privatization}
Per-block private histogram in shared; merge once at end:
\begin{lstlisting}
__shared__ unsigned histo_p[7];
if(threadIdx.x<7) histo_p[threadIdx.x]=0;
__syncthreads();
while(i<size){
  atomicAdd(&histo_p[(buf[i]-'a')/4],1);
  i+=stride; }
__syncthreads();
if(threadIdx.x<7)
  atomicAdd(&histo[threadIdx.x],
            histo_p[threadIdx.x]);
\end{lstlisting}
Requires op \textbf{associative + commutative} and copy fits in shared. Too big? Partial privatization (hot bins only).

\sect{L12: Multi-GPU}
\sub{Device management}
CUDA calls target the \emph{current} device; \texttt{cudaSetDevice(i)} switches. Switching while async work is in flight is fine, both GPUs run concurrently:
\begin{lstlisting}
cudaSetDevice(0); k<<<...>>>();
cudaSetDevice(1); k<<<...>>>(); //0 still runs
\end{lstlisting}

\sub{P2P}
\begin{lstlisting}
cudaDeviceCanAccessPeer(&can,dX,dY);
cudaDeviceEnablePeerAccess(peer,0);
cudaMemcpyPeerAsync(dst,dDev,src,sDev,n,s);
\end{lstlisting}
With P2P: shortest PCIe path, no host staging. Without: routed through host.\\
\textbf{Numbers (PCIe gen2):} P2P $\sim$6.6 GB/s simplex, $\sim$12 duplex (vs 5 via host); gen3 doubles. 4 GPUs via switches $\sim$15 GB/s aggregate.

\sub{Halo exchange (1D decomp)}
Stage 1: all send right, recv left. Stage 2: all send left, recv right. Concurrent copies, no link contention.

\sub{NUMA}
Dual-IOH: remote D2H crosses QPI, extra hop. \textbf{Local $\sim$6.3 GB/s vs remote $\sim$4.3 GB/s.} Dual-IOH \textbf{blocks P2P across IOHs} (falls back to host staging). Pin threads near GPU: \texttt{numactl}, \texttt{GOMP\_CPU\_AFFINITY}.

\sect{L13: GPU Sharing, MPS}
\sub{Problem + idea}
Strong-scaling MPI: more ranks = less GPU work each; each process has its own context and the GPU \textbf{serializes contexts} (switch overhead, no overlap). \textbf{MPS} = one shared context via a server process; kernels/copies from different processes can overlap.

\sub{Requirements}
Linux only; UVA; Tesla CC $\geq$ 3.5; CUDA 5.5+. Exclusive mode applies to the MPS \emph{server}, not clients.

\sub{Hyper-Q / concurrency}
Concurrent kernels since Fermi (CC 2.0) on different streams if resources allow: Fermi 16, Kepler 32. \textbf{Kepler Hyper-Q:} 32 HW work queues, one per stream, kills false inter-stream dependencies. Under MPS each client gets 2 queues per channel, 32 total. Streams/client $\leq$ channels: fine. Streams/client $>$ channels: \textbf{false dependency, serialization}.

\sub{Setup (no code changes)}
\begin{lstlisting}
export CUDA_VISIBLE_DEVICES=0
nvidia-smi -i 0 -c EXCLUSIVE_PROCESS
nvidia-cuda-mps-control -d
\end{lstlisting}
Daemon spawns \texttt{nvidia-cuda-mps-server}; later processes route CUDA work through it. MPS does \textbf{not} spread work across GPUs (app handles affinity). Profile: \texttt{nvprof --profile-all-processes -o f\_\%p}.

\sect{L14: Dynamic Parallelism}
\sub{Core idea}
Kepler+ (CC 3.5+): a kernel can launch child kernels with normal \texttt{<<<g,b>>>} syntax. Unlocks: device-side library calls (cuBLAS), data-dependent parallelism (irregular bounds \texttt{for j=1..x[i]}), recursion, adaptive mesh refinement, batched nested parallelism (CPU freed).

\sub{Rules (memorize)}
\begin{itemize}
\item \textbf{Launch is per-thread}: 2 threads at the launch line $\Rightarrow$ 2 child grids.
\item \textbf{Sync is per-block}: device \texttt{cudaDeviceSynchronize()} waits for launches by \emph{any} thread in the block.
\item \texttt{cudaDeviceSynchronize()} does \textbf{NOT} imply \texttt{\_\_syncthreads()}; add it if you need a block barrier after.
\item All device-side launches and copies are \textbf{async}; no sync launch.
\item Streams/events are per-block, cannot pass to children. Constants set from host; textures/surfaces host-bound; ECC errors at host.
\end{itemize}

\sub{Pattern}
\begin{lstlisting}
__global__ void dyn(float *d){
  int t=threadIdx.x;
  if(t%2) buf[t/2]=d[t]+d[t+1];
  __syncthreads();
  if(t==0){
    child<<<128,256>>>(buf);
    cudaDeviceSynchronize(); }
  __syncthreads(); // still required!
  cudaMemcpyAsync(d,buf,1024);
  cudaDeviceSynchronize();
}
\end{lstlisting}
Library call: one thread calls \texttt{cublasDgemm} + \texttt{cudaDeviceSynchronize()}, then \texttt{\_\_syncthreads()} so all threads see the result.

\sect{Carryover from Midterm 1}
Warp = \textbf{32 threads}. Warps/block $= \lceil \text{blockDim}/32 \rceil$.\\
2D grid for $W{\times}H$ image, block $B_x{\times}B_y$: grid $= (\lceil W/B_x\rceil, \lceil H/B_y\rceil)$; ceiling div: \texttt{(n+b-1)/b}.\\
Divergent warps: boundary blocks where the \texttt{if(Row<H\&\&Col<W)} test splits a warp.\\
Coalescing check: is the index \texttt{base + threadIdx.x}?

\sect{Quick numbers \& trivia}
\begin{itemize}
\item atomic returns \textbf{old} value.
\item \texttt{cudaMemAdvise} = hints, no movement; \texttt{cudaMemPrefetchAsync} = moves.
\item Atomic latency: DRAM $\gg$ L2 ($\sim$10x faster) $\gg$ shared ($\sim$100x faster).
\item Tiled SGEMM: FLOPs/load $=T$; shared/blk $=8T^2$ bytes (2 float tiles).
\item P2P gen2: 6.6 simplex / 12 duplex GB/s; host path 5.
\item NUMA: 6.3 local vs 4.3 remote GB/s D2H.
\item Hyper-Q: 32 queues; MPS client: 2 per channel.
\item DP needs CC 3.5+; MPS needs CC 3.5+, Linux, UVA.
\item Pre-Pascal UM: CPU touch during kernel segfaults; Pascal+ page-faults.
\end{itemize}

\sect{Worked example: launch math}
$1000{\times}1000$ image, $16{\times}16$ blocks:\\
Grid $= \lceil 1000/16\rceil^2 = 63{\times}63 = 3969$ blocks.\\
Threads $= 3969{\cdot}256 = 1{,}016{,}064$ launched ($16{,}064$ idle past bounds).\\
Warps/block $= 256/32 = 8$; total warps $= 31{,}752$.\\
Divergence: a $16{\times}16$ block has warps spanning 2 rows of 16. Edge blocks (right col, bottom row) have warps where \texttt{if(Row<H\&\&Col<W)} splits inside the warp $\Rightarrow$ divergent; interior blocks have none.

\sub{Tiled occupancy check}
T=32, 48KB shared/SM, 2048 thr/SM max: shared allows $48/8=6$ blocks but thread limit allows $2048/1024=2$ $\Rightarrow$ \textbf{2 blocks/SM} (thread bound). Always take the min over all limits (threads, shared, registers, blocks).

\sect{API quick reference}
\begin{lstlisting}
kernel<<<grid,block,shmemB,stream>>>(...);
dim3 g(gx,gy), b(bx,by);
cudaMalloc(&p,N); cudaFree(p);
cudaMemcpy(dst,src,N,cudaMemcpyHostToDevice);
cudaMemcpyAsync(d,s,N,kind,stream);
cudaStreamCreate(&s); cudaStreamDestroy(s);
cudaMallocManaged(&p,N);
cudaMemPrefetchAsync(p,N,dev,stream);
cudaMemAdvise(p,N,hint,dev);
cudaSetDevice(i); cudaGetDeviceCount(&n);
cudaDeviceCanAccessPeer(&ok,a,b);
cudaDeviceEnablePeerAccess(peer,0);
cudaMemcpyPeerAsync(d,dDev,s,sDev,N,st);
cudaDeviceSynchronize();
old = atomicAdd(&x,v);
\end{lstlisting}

\sect{Which optimization when}
\begin{itemize}
\item Redundant global reads across threads $\Rightarrow$ \textbf{tiling} in shared memory.
\item Strided / scattered warp access $\Rightarrow$ restructure for \textbf{coalescing} (or stage through shared).
\item Many threads update few locations $\Rightarrow$ \textbf{atomics + privatization}.
\item Data too big for one GPU or need more FLOPs $\Rightarrow$ \textbf{multi-GPU + P2P} halo exchange.
\item Many small processes underusing GPU $\Rightarrow$ \textbf{MPS}.
\item Work amount unknown until runtime $\Rightarrow$ \textbf{dynamic parallelism}.
\item Complex pointer structures, prototyping $\Rightarrow$ \textbf{unified memory} (+prefetch/advise to claw back perf).
\end{itemize}

\end{multicols*}
\end{document}
